\section{背景术语速查}

正文的词条在「对应」一段里反复借用一批数学名词：相干层、导出范畴、Dolbeault 上同调几乎贯穿 B-模型、B-膜与几何 Langlands 的每一条，Nahm 方程与 Nahm 定理则夹在单极子和维数约化的讨论中。它们都是被当作已知语言写下的，KW06 原文也只引用、不展开；附录 A 的总表更点明，其中八条属于「原文只引用而未实质使用」。读者若只读本词典，这些词就停在字面上——知道名字，却不知道它指什么、为什么出现在这里。本附录为这八条各给一个最小可用的定义：每条按「一句话、在 KW06 里怎么用、最小定义、易错点、延伸」五项写，控制在几行以内，只求把正文的对应句读通，不展开证明，也不替代任何教材。

八条的顺序沿用附录 A：先 Nahm 方程与 Nahm 定理，再相干层与导出范畴，然后环路群与 opers，最后 Dolbeault 上同调与 Kac--Moody 代数。「在 KW06 里怎么用」一栏凡写「原文未出现该词」的，都是核对原文后的结论，不是猜测；这类术语的价值在于说明原文某处背后的数学语言，而不是原文用过的行话。

\subsection{Nahm 方程（Nahm equations）}

\textbf{一句话}：描述单极子模空间的一组矩阵常微分方程。

\textbf{在 KW06 里怎么用}：原文没有写出这组方程。§2.1 只借「Nahm 定理」之名说明十维的来源，单极子本身在 §6.2、§9--§10 由 Bogomolny 方程处理；Nahm 方程是同一模空间的另一套（ADHMN）描述，本词典在「单极子」条以 Nahm 数据、Nahm 极点之名用到它。

\textbf{最小定义}：固定磁荷 $k$，未知函数是区间上三个取值于 $\mathfrak{u}(k)$ 的矩阵 $T_1,T_2,T_3$，满足 $\frac{dT_i}{ds}=\frac12\epsilon_{ijk}[T_j,T_k]$；端点允许单极点，留数给出单极子的位置等参数。ADHMN 构造给出它与单极子模空间的对应。

\textbf{易错点}：Nahm 方程不是 Hitchin 方程：前者是区间上的矩阵常微分方程，后者是黎曼面上的全纯方程；它也与 §2.1 的 Nahm 定理无关。

\textbf{延伸}：Gaiotto--Witten 用 Nahm 极点刻画超对称边界条件，Frenkel 2009 综述有简要回顾。

\subsection{Nahm 定理（Nahm's theorem）}

\textbf{一句话}：无引力的超对称理论中，含规范场的超对称杨--米尔斯理论最高只能存在于十维。

\textbf{在 KW06 里怎么用}：§2.1（p.6）引它解释为什么 $\mathcal N=4$ 杨--米尔斯最自然地来自十维理论的维数约化，这句引用把 §2 的超对称回顾与 §3 的扭化接了起来。译本在此补了一条译注，复述 1978 年 Nahm 的原始论证。

\textbf{最小定义}：超对称代数中，超荷的旋量表示维数随维数增长。十维的 Majorana--Weyl 旋量是十六分量，恰好容得下规范多重态；十一维最小的旋量已是三十二分量，超对称杨--米尔斯不再存在，只允许超引力，所以杨--米尔斯的上界是十维。

\textbf{易错点}：定理约束的是超对称杨--米尔斯，而不是所有超对称理论：十一维超引力就是现成的反例，两条上界恰好差一维。

\textbf{延伸}：原文 §2.1 与译本译注。

\subsection{相干层（coherent sheaf）}

\textbf{一句话}：全纯函数层上的模，且局部有限表现（有限个生成元与有限个关系）。

\textbf{在 KW06 里怎么用}：§3.1 的扭化族在 §3.2 中于 $t=\pm i$ 给出 B-模型，B-模型的对象层面就是相干层；§5.6 讨论 $\mathcal M_H$ 上的 $(B,B,B)$ 膜时，说它是「填满空间、带全纯 Chan--Paton 丛的膜」，丛与它承载的层即相干层。B-膜的范畴因此写成 $D^b(\mathrm{Coh}(X))$（见下条）。

\textbf{最小定义}：设 $X$ 是复流形，$\mathcal O_X$ 为全纯函数层。一个 $\mathcal O_X$-模 $\mathcal F$ 称为相干层，若它在每点附近有有限表现：存在正合列 $\mathcal O_U^{\oplus p}\to\mathcal O_U^{\oplus q}\to\mathcal F|_U\to0$。等价地，$\mathcal F$ 局部由有限个截面生成，且任意有限个截面之间的关系子层也有限生成；仿射情形下这正是有限生成模（Serre）。

\textbf{易错点}：相干不等于局部自由：全纯向量丛只是局部自由的特例，带挠的层可以相干却不局部自由。

\textbf{延伸}：B-膜的态空间是 $\mathrm{Ext}^\bullet$，可用 Dolbeault 复形计算（见本附录 Dolbeault 上同调条）。

\subsection{导出范畴（derived category）}

\textbf{一句话}：以复形为对象、把拟同构当作同构的范畴；它记住上同调与高阶态射。

\textbf{在 KW06 里怎么用}：§3.1 说扭化理论中规范不变的算子与态构成复形，按鬼数分次、以 $Q$ 为微分，取 $Q$-上同调得拓扑不变量；导出范畴是这套语言的范畴化，B-膜的范畴即 $D^b(\mathrm{Coh}(X))$。原文的 Fourier--Mukai 与 §11 的膜--模对应只在启发式层面使用它。

\textbf{最小定义}：设 $\mathcal A$ 是 abelian 范畴。$D(\mathcal A)$ 的对象是 $\mathcal A$ 中的复形，态射由链映射「把拟同构形式地变成同构」得到（Verdier 局部化）；$D^b(\mathcal A)$ 只取有界复形。两个复形若拟同构，在导出范畴中就是同构对象，但复形的信息并未丢失。

\textbf{易错点}：导出范畴不是复形对拟同构的商，也不是同伦范畴：它由同伦范畴再对拟同构作局部化得到，简单地把拟同构「去掉」会丢掉高阶信息。

\textbf{延伸}：$D^b(\mathrm{Coh}(X))$ 上的 Fourier--Mukai 以核的复形给出（Huybrechts 的综述）。

\subsection{环路群（loop group）}

\textbf{一句话}：取值于形式 Laurent 级数域的群，即 $G(\!(z)\!)$。

\textbf{在 KW06 里怎么用}：§9.3 定义仿射 Grassmann 流形 $\mathrm{Gr}_N=GL(N,\mathbb C(\!(z)\!))/GL(N,\mathbb C[[z]])$，并注明它就是圈 Grassmann 流形；原文用它的 Schubert 簇描述 't Hooft 算子 Hecke 修正空间的胞腔分解，但随即说明「本文其实不会真正使用这个框架」。环路群本身没有单独出场。

\textbf{最小定义}：设 $G$ 是复线性代数群。$G(\!(z)\!)$ 是系数在形式 Laurent 级数 $\mathbb C(\!(z)\!)$ 中、作为矩阵可逆的元素组成的群；$G[[z]]$ 是正幂级数子群，商 $G(\!(z)\!)/G[[z]]$ 即仿射 Grassmann 流形。其李代数是环路代数 $\mathfrak g\otimes\mathbb C(\!(z)\!)$，加中心扩张后成为仿射 Kac--Moody 代数（见末条）。

\textbf{易错点}：$G(\!(z)\!)$ 是形式环路群，不是拓扑环路空间 $\mathrm{Map}(S^1,G)$：它只取 Laurent 级数、不问收敛，因而是代数对象。

\textbf{延伸}：Pressley--Segal《Loop Groups》是标准参考。

\subsection{opers}

\textbf{一句话}：一种带额外结构的平坦联络：在完整旗上，联络的逐次商映射全为同构。

\textbf{在 KW06 里怎么用}：原文没有出现这个词；§12 讨论膜时出现的平坦丛是 $G_{\mathbb C}$ 上的（§12.4），没有加旗结构。本词典按几何 Langlands 的约定把 oper 记在对偶群一侧，写成 ${}^L G$-opers：Langlands 参数是 ${}^L G$ 上的平坦丛，旗结构加在它上面；若加在 $G$ 一侧则是 $G$-opers，二者一般不同（$G=SL_2$ 时 ${}^L G=PGL_2$）。

\textbf{最小定义}：设 $G$ 单连通单群、$B$ 为 Borel 子群。$X$ 上的 $G$-oper 是三元组 $(E,\nabla,E_B)$：$G$-丛 $E$、平坦联络 $\nabla$、Borel 约化 $E_B$（即完整旗 $0\subset F^1\subset\cdots\subset E$），满足 $\nabla(F^i)\subset F^{i-1}\otimes\Omega_X^1$，且逐次商间的映射 $F^i/F^{i+1}\to F^{i-1}/F^i\otimes\Omega_X^1$ 都是同构。

\textbf{易错点}：别把 ${}^L G$-opers 与 $G$-opers 互换：Langlands 参数取在对偶群上；也别把 oper 当成一般平坦联络，旗结构才是额外条件，它使 oper 的模空间只有平坦丛模空间的一半维数。

\subsection{Dolbeault 上同调（Dolbeault cohomology）}

\textbf{一句话}：复流形上用 $\bar\partial$ 算子算出的上同调，记作 $H^{p,q}_{\bar\partial}(X)$。

\textbf{在 KW06 里怎么用}：§5.4 说 B-模型的手征环是双分次环 $H^q(X,\wedge^pT_X)$，$p=q=0$ 的子环是全纯函数环；这个上同调就是取值于全纯丛 $\wedge^pT_X$ 的 Dolbeault 上同调。原文直接写层上同调，不写 $\bar\partial$ 记号，B-膜的态空间 $\mathrm{Ext}^\bullet$ 则用 Dolbeault 复形计算。

\textbf{最小定义}：设 $E$ 是 $X$ 上的全纯向量丛，$\Omega^{p,q}(X,E)$ 为取值于 $E$ 的 $(p,q)$ 型形式。$\bar\partial_E$ 把 $(p,q)$ 型映到 $(p,q+1)$ 型且平方为零，定义 $H^{p,q}_{\bar\partial}(X,E)=\ker\bar\partial_E/\operatorname{im}\bar\partial_E$；Dolbeault 定理说它同构于 $H^q(X,\Omega^p\otimes E)$，$E$ 平凡时即 $H^q(X,\Omega^p)$。

\textbf{易错点}：Dolbeault 上同调不是 de Rham 上同调：$\bar\partial$ 保持双型，$d$ 却混合各型；上标 $p,q$ 是形式类型，别与层上同调的次数混淆。

\textbf{延伸}：紧 Kähler 流形上 Hodge 理论给出它与调和形式的对应。

\subsection{Kac--Moody 代数（Kac--Moody algebra）}

\textbf{一句话}：由广义 Cartan 矩阵按 Serre 关系生成的李代数；仿射型正是带中心扩张的环路代数，即仿射李代数。

\textbf{在 KW06 里怎么用}：原文没有出现这个词；§9.3 的环路群 $GL(N,\mathbb C(\!(z)\!))$ 的李代数作中心扩张就是仿射 Kac--Moody 代数。几何 Langlands 的 Beilinson--Drinfeld 构造在临界层用它构造 opers 对应的 Hecke 特征层；KW06 走 §11 的 A-膜与 D-模路线，因此不提。

\textbf{最小定义}：设 $\mathfrak g$ 是单李代数，$A$ 为其仿射型广义 Cartan 矩阵；Kac--Moody 代数 $\mathfrak g(A)$ 由 $e_i,f_i,h_i$ 及 Serre 关系生成。仿射型可用环路代数实现：$\hat{\mathfrak g}=\mathfrak g\otimes\mathbb C[z,z^{-1}]\oplus\mathbb C K\oplus\mathbb C d$，$K$ 为中心元且 $[K,\cdot]=0$，$[x\otimes z^m,y\otimes z^n]=[x,y]\otimes z^{m+n}+m\delta_{m+n,0}(x,y)K$；$K$ 的特征值称水平（level），临界层为 $k=-h^\vee$（$h^\vee$ 是对偶 Coxeter 数）。

\textbf{易错点}：Kac--Moody 代数不全是仿射型（还有有限型与不定型）；仿射型也不同于无中心扩张的环路代数，差的正是中心 $K$ 与导子 $d$，水平与临界层的概念都依赖 $K$。

\textbf{延伸}：Kac《Infinite Dimensional Lie Algebras》；Frenkel--Ben-Zvi《Vertex Algebras and Algebraic Curves》。
